\chapter{\IfLanguageName{dutch}{Analyse}{analysis}}%
\label{ch:analyse}

\section{\IfLanguageName{dutch}{Architecturale Analyse}{Architectural analysis}}%
\label{sec:architecturale-analyse}

In dit hoofdstuk worden verschillende mainframe opstellingen vergeleken en hun impact op RPO, RTO en dataconsistentie geanalyseerd.

\subsection{\IfLanguageName{dutch}{Twee Synchrone Sites}{Two Synchronous Sites}}%
\label{sec:twee-synchrone-sites}

Een veelgebruikte opstelling is een configuratie met twee sites die gesynchroniseerd zijn met GDPS Metro Mirror.
GDPS Metro Mirror tussen twee sites is een veelvoorkomende opstelling bij omgevingen waarin dataverlies niet acceptabel is, zoals financiële transactieverwerkingssystemen \autocite{ivanov2023unveiling}.
In deze sectie wordt deze technologie nog eens uitgelegd.
De beperkingen van de technologie worden besproken, samen met de mogelijke faalscenario's en risico's die hiermee gepaard gaan.
Alle bevindingen worden kort samengevat in een tabel op het einde van deze sectie.

\subsubsection{\IfLanguageName{dutch}{Werking van GDPS Metro Mirror}{Operation of GDPS Metro Mirror}}%
\label{sec:twee-synchrone-sites-werking-metro-mirror}

GDPS Metro Mirror is een synchrone replicatietechnologie waarbij schrijfoperaties pas als voltooid worden beschouwd wanneer ze zowel op de primaire als de secundaire site zijn bevestigd.
Dit zorgt voor een RPO van nul, wat betekent dat er geen dataverlies optreedt bij een uitval van de primaire site.
Deze eigenschap maakt Metro Mirror bijzonder geschikt voor kritieke systemen waar dataintegriteit essentieel is, zoals in bankomgevingen \autocite{gomes_2025_redbooks}.

\subsubsection{\IfLanguageName{dutch}{Afstandsbeperking en Impact op Prestaties}{Length Limit and Impact on Performance}}%
\label{sec:twee-synchrone-sites-beperkingen}

Een belangrijke beperking van GDPS Metro Mirror is de maximale afstand tussen de twee sites.
IBM specifieert dat Metro Mirror een maximale afstand toelaat van 300 km \autocite{gomes_2025_redbooks}.
Deze beperking is het gevolg van de synchrone aard van de replicatie.
Elke schrijfoperatie moet namelijk worden bevestigd door de secundaire site voordat de primaire site de operatie als voltooid beschouwt.
Hierdoor neemt de responstijd van applicaties toe naarmate de afstand groter wordt, omdat de vertraging op netwerkniveau toeneemt.
De RTO van GDPS Metro Mirror varieert van enkele seconden tot een uur, afhankelijk van de mate van automatisering. 
Met geautomatiseerde uitvalbeveiliging zoals GDPS HyperSwap kan de RTO worden teruggebracht tot enkele seconden, omdat de secundaire site onmiddellijk kan worden geactiveerd. 
Zonder automatisering kan de RTO echter oplopen tot een uur, omdat handmatige interventie nodig is om de secundaire site operationeel te maken \autocite{li_2019_planet_mainframe_gdps}.
\subsubsection{\IfLanguageName{dutch}{Faalscenario's en Risico's}{Possible Failures and Risks}}
\label{sec:twee-synchrone-sites-failures-risks}

Hoewel Metro Mirror een RPO van nul biedt, zijn er risico's verbonden aan deze opstelling.
Een mogelijk risico is regionale uitval. 
Omdat GDPS Metro Mirror beperkt is tot maximaal 300 km kunnen regionale incidenten zoals stroomuitval, natuurrampen en cyberaanvallen beide sites treffen.
In dit geval is er geen beschikbare secundaire site om de dienstverlening voort te zetten \autocite{li_2019_planet_mainframe_gdps}.

\begin{table}[H]
\begin{tabular}{|l|l|}
\hline
Replicatietype        & Synchroon                               \\ \hline
RPO                  & nul                                       \\ \hline
RTO                  & Seconden tot een uur                      \\ \hline
Maximale afstand     & $\sim$300 km                            \\ \hline
Impact op prestaties & Toename responstijd bij grotere afstand \\ \hline
Faalscenario's       & Geregionaliseerde uitval               \\ \hline
\end{tabular}
\caption{Kenmerken van een opstelling met twee synchrone sites (GDPS Metro Mirror).}
\end{table}

\begin{figure}[H]
    \centering
    \includegraphics[width=0.8\textwidth]{ibm_metro_mirror.png}
    \caption{Metro Mirror: synchronous remote copy. Bron: \cite{ivanov2023unveiling}.}
    \label{fig:ibm_metro_mirror}
\end{figure}

\subsection{\IfLanguageName{dutch}{Twee Asynchrone Sites}{Two Asynchronous Sites}}%
\label{sec:twee-asynchrone-sites}

Een andere voorkomende opstelling is een configuratie van twee sites die asynchroon verbonden zijn doormiddel van GDPS Global Mirror.
Het laat toe om sites met elkaar te verbinden die meer dan 200 km van elkaar verwijderd zijn \autocite{gomes_2025_redbooks}.
In deze sectie wordt onderzocht hoe deze asynchrone opstelling zich verhoudt tot RPO en RTO.
Concreet wordt de werking herhaald en worden de voordelen en beperkingen van deze opstelling opgelijst.
Mogelijkse faalscenario's en risico's worden besproken.
Al deze informatie wordt, zoals bij de vorige opstelling, finaal samengevat in een tabel.

\subsubsection{\IfLanguageName{dutch}{Asynchrone replicatie}{Asynchronous replication}}%
\label{sec:asynchrone-replicatie}

Bij asynchrone replicatie worden schrijfoperaties op de primaire site onmiddelijk bevestigd aan de applicatie, zonder te wachten op bevestiging van de externe site.
Dit zorgt ervoor dat er minimale impact is op de prestaties van de primaire site, maar introduceert een replicatievertraging.
Deze vertraging bepaalt de RPO, want hoe groter de vertraging, hoe meer data er verloren kan gaan bij het falen van de primaire site.

\subsubsection{\IfLanguageName{dutch}{Werking van GDPS Global Mirror}{Operation of GDPS Global Mirror}}%
\label{sec:twee-asynchrone-sites-werking-global-mirror}

GDPS Global Mirror is een replicatietechnologie die toelaat om op asynchrone wijze data te kopiëren tussen sites \autocite{white_2020_redbooks}.
Door de asynchrone werking van Global Mirror, zijn er geen limitaties op de afstand tussen de sites.
Het maakt gebruik van consistency groups om periodiek een consistent herstelpunt te creëren \autocite{ibm_2023_ibm_global_mirror}.
Dit zorgt voor een RPO van seconden tot een minuut, afhankelijk van de frequentie van de consistency groups en de beschikbare bandbreedte \autocite{raften_2011_gdps_executive_overview}.
De RTO varieert van 30 seconden tot 2 uur, afhankelijk van de mate van automatisering en de complexiteit van het herstelproces.

\subsubsection{\IfLanguageName{dutch}{Voordelen en beperkingen}{Benefits and limits}}%
\label{sec:twee-asynchrone-sites-voordelen-beperkingen}

Asynchrone replicatie doormiddel van GDPS Global Mirror biedt verschillende voordelen.
Omdat de primaire site niet moet wachten op bevestiging van de secundaire site, is er geen invloed op de prestaties van applicaties \autocite{gomes_2025_redbooks}.
Deze technologie kan tevens worden toegepast over lange afstanden, waardoor de impact van regionale incidenten kan worden geminimaliseerd.
Echter zijn er ook beperkingen, want Global Mirror biedt geen garantie bij het volledig voorkomen van dataverlies \autocite{gomes_2025_redbooks}.
Consistency groups kunnen achterlopen en niet volledig real-time consistentie garanderen.
Het herstelproces is ook complex doordat er sprake kan zijn van inconsistenties.
Ook is voldoende bandbreedte en minimale vertraging op netwerkniveau noodzakelijk om de effectiviteit van deze asynchrone replicatietechnologie te waarborgen.

\subsubsection{\IfLanguageName{dutch}{Faalscenario's en Risico's}{Possible Failures and Risks}}
\label{sec:twee-asynchrone-sites-failures-risks}

Asynchrone replicatie met GDPS Global Mirror introduceert specifieke faalscenario's en risico's, met name gerelateerd aan de replicatievertraging en de asynchrone aard van de technologie.

Bij asynchrone replicatie met GDPS Global Mirror onstaat een vertraging tussen het moment waarop de data werd geschreven op de primaire site en het moment waarop dezelfde data beschikbaar is op de secundaire site.
Deze vertraging is onlosmakelijk verbonden aan de asynchrone werking en wordt beïnvloed door verschillende factoren.
Een van deze factoren is de voorheen besproken frequentie van consistency groups. 
Hoe vaker een consistency group wordt gemaakt, hoe minder de discrepantie tussen de data op de primaire site en de data op de secundaire site.
Bandbreedte kan ook vertragingen veroorzaken doordat het langer duurt tot de data gerepliceerd is op de secundaire site.
Hierbij kan de replicatievertraging geminimaliseerd worden door zaken zoals compressietechnieken. 
Het implementeren van geautomatiseerde mechanismen zoals GDPS HyperSwap kan er voor zorgen dat schakelen naar de secundaire versneld wordt, waardoor de RTO ook kan worden verkort.

\begin{table}[H]
\begin{tabular}{|l|l|}
\hline
Replicatietype       & Asynchroon                                 \\ \hline
RPO                  & Seconden tot een minuut                      \\ \hline
RTO                  & 30 seconden tot 2 uur                      \\ \hline
Maximale afstand     & Wereldwijd                                 \\ \hline
Impact op prestaties & Minimaal (geen wachttijd)                  \\ \hline
Faalscenario's       & Replicatievertraging                       \\ \hline
\end{tabular}
\caption{Kenmerken van een opstelling met twee asynchrone sites (GDPS Global Mirror).}
\end{table}

\subsection{\IfLanguageName{dutch}{Twee Synchrone Sites met een Asynchrone site}{Two Synchronous Sites with a Asynchronous Site}}%
\label{sec:hybride-opstelling} 

Een hybride replicatie-opstelling combineert de voordelen van synchrone en asynchrone replicatie door twee synchrone sites te koppelen aan een derde asynchrone site.
In deze sectie zal worden onderzocht of een hybride opstelling de nadelen van Metro Mirror en Global Mirror kan minimaliseren door de twee technologieën te combineren.
Er wordt onderzocht of het probleem kan worden opgelost bij Metro Mirror waarbij er geen bescherming is voor regionale uitval.
Tevens wordt ook onderzocht of door deze hybride opstelling de data inconsistentie bij een asynchrone opstelling met Global Mirror kan worden vermeden.

\subsubsection{\IfLanguageName{dutch}{Werking Hybride Opstelling}{Operation of Hybrid Setup}}
\label{sec:werking-hybride-opstelling}

In een hybride opstelling worden twee sites gesynchroniseerd doormiddel van GDPS Metro Mirror (bijvoorbeeld site A en site B), en wordt een derde site (site C) asynchroon gerepliceerd met GDPS Metro Mirror.
Hierbij geldt dat site A en site B gesynchroniseerd zijn met een RPO van 0 binnen de maximale afstand waarbij Metro Mirror toegelaten is.
Zoals eerder besproken wordt tussen deze twee sites met Metro Mirror, elke schrijfoperaties pas bevestigd wanneer beide sites de data hebben ontvangen.
Dit is noodzakelijk om dataverlies te vermijden wanneer een van de twee sites uitvalt. 
De RTO varieert van seconden tot een uur \autocite{li_2019_planet_mainframe_gdps}.
Dit hangt af van de mate van automatisering.
Om de hybride opstelling compleet te maken, is er een aanvullende asynchrone site die zich buiten de maximale afstand van GDPS Metro Mirror bevindt.
Een van de twee synchrone sites (site A of site B) repliceert hierbij op asynchrone wijze naar site C met behulp van GDPS Global Mirror.
Tevens, zoals eerder besproken, biedt deze manier van replicatie een RPO van seconden tot een minuut.
De RTO varieert van 30 seconden tot twee uur.
Dit hangt af van de frequentie van consistency groups en de beschikbare bandbreedte \autocite{gomes_2025_redbooks}.
De asynchrone site (site C in dit geval) kan wereldwijd worden geplaatst waardoor er bescherming geboden wordt tegen regionale incidenten zoals natuurrampen en stroomuitval.

\subsubsection{\IfLanguageName{dutch}{Voordelen Hybride Opstelling}{Pros of Hybrid Setup}}
\label{sec:voordelen-hybride-opstelling}

Het gebruik van een hybride opstelling waarbij zowel een synchrone replicatietechnologie (GDPS Metro Mirror) als een asynchrone replicatietechnologie (GDPS Global Mirror) wordt gebruikt, biedt voordelen.
De synchrone sites, site A en B, garanderen een RPO van 0, wat cruciaal is voor systemen zoals financiële transactieverwerking.
De asynchrone site, site C, biedt bescherming tegen regionale uitval door data buiten het maximale GDPS Metro Mirror bereik van de synchrone sites te plaatsen.
Door de voordelen van synchrone replicatie te combineren met die van asynchrone replicatie, is er flexibiliteit en schaalbaarheid \autocite{gomes_2025_redbooks}.

\subsubsection{\IfLanguageName{dutch}{Faalscenario's en Risico's}{Possible Failures and Risks}}
\label{sec:faalscenarios-risicos-hybride-opstelling}

Hoewel de hybride opstelling veel voordelen biedt, introduceert het ook specifieke faalscenario's en risico's.
Als site A én site B uitvallen door regionale uitval moet er worden overgeschakeld naar site C om geen verlies van kritieke diensten te veroorzaken.
Bij het overschakelen naar site C kan de RPO beperkt worden tot seconden.
Echter, door de asynchrone werking van site C, kan de RTO lopen van 30 seconden tot twee uur.
Om de impact op kritieke diensten zoveel mogelijk te mitigeren, is het noodzakelijk om de overschakeling naar site C te automatiseren.
Frequentere consistency groups zijn hierbij ook nuttig.
De vertraging tussen site B en C bij het asynchroon repliceren kan ook problemen veroorzaken als er netwerkproblemen.
Om frequent consistency groups te kunnen bewaren op site C, moet de netwerkverbinding tussen site B en C stabiel zijn.
Bovendien is het mogelijk dat er data verloren gaat wanneer site B uitvalt.
Dit heeft dan ook een impact op de RPO voor site C.
Als er inconsistenties onstaan in site C, dan kan het zijn dat handmatige interventie nodig waarbij ook de RTO stijgt.
Technologieën zoals GDPS HyperSwap kunnen hierbij helpen om geautomatiseerde failover te creëren.
Echter moet ook worden nagedacht over consistency groups en dat site C pas operationeel kan worden zodra alle consistency groups zijn gerepliceerd.
Om site C operationeel te maken moeten ook alle schrijfoperaties van de applicaties op de primaire site gepauzeerd worden totdat alle updates verstuurd zijn naar de herstelsite (in dit geval site C) \autocite{ibm_2021_ibm_copy_consistency_groups}.

\subsubsection{\IfLanguageName{dutch}{Impact op RPO en RTO}{Impact on RPO and RTO}}
\label{sec:hybride-opstelling-impact-rpo-rto}

\begin{table}[H]
\begin{tabularx}{\textwidth}{|p{3cm}|p{2cm}|p{2cm}|X|}
\hline
Site & RPO  & RTO & Opmerking \\ \hline
Synchrone sites (A en B) & 0 & Seconden tot een uur & Biedt een lage RPO en snelle failover binnen regionale omgeving.  \\ \hline
Asynchrone site (B repliceert naar C) & Seconden tot een minuut & 30 seconden tot twee uur & Biedt bescherming tegen regionale uitval, maar heeft wel een hogere RPO en RTO door de asynchrone werking. \\ \hline
\end{tabularx}
\caption{Kenmerken van een hybride opstelling (2 sites gesynchroniseerd met GDPS Metro Mirror samen met een extra site gesynchroniseerd met GDPS Global Mirror).}
\end{table}

\subsubsection{\IfLanguageName{dutch}{Mitigatiestrategieën}{Mitigation Strategies}}
\label{sec:hybride-opstelling-mitigatiestrategie}

Om de risico's die gepaard gaan met een hybride opstelling te vermijden, kunnen de volgende strategieën worden toegepast.
Wanneer logs in Db2 voor z/OS worden gerepliceerd naar een externe site is het nuttig om het meest performante opslagsysteem ter beschikking te stellen, omdat de snelheid van het schrijven van log data vanuit de log buffers naar de active log ook te maken heeft met fysieke hardware die de schrijfoperaties doet.
Ook is het belangrijk om te weten dat kleine log buffers ervoor zorgen dat er meer schrijfoperaties noodzakelijk zijn.
Het is daarom nuttig om de log buffers van zoveel mogelijk plaats in de virtuele opslag te voorzien, zodat het aantal schrijfoperaties kunnen worden verminderd.
Schrijfoperaties naar fysieke opslagsystemen ten opzichte van schrijven naar virtuele opslag is namelijk trager \autocite{wisniewski_2017_benefits_more_memory}.
De boodschap is om paging (het uitwisselen van gegevens tussen geheugen en fysieke opslag) zoveel mogelijk te vermijden.
Het deel van het z/OS besturingsssysteem die zich ontfermt over het beheer van virtuele opslag is ook zo ontworpen om de CPU-prestaties aanzienlijk te verbeteren, doordat paging zoveel mogelijk wordt vermeden \autocite{ibm_2026_zos_virtual_storage_management}.

Naast efficiënt beheer van log buffers, zijn er uiteraard nog andere strategieën om de impact van risico's die gepaard gaan met een hybride opstelling te verminderen.
Zo kan automatische failover worden geïmplementeerd met GDPS HyperSwap tussen site A en B.
Bij een ongeplande HyperSwap detecteert het systeem automatisch een storing in volumes en wordt er overschakeld naar een andere set volumes.
De HyperSwap functionaliteit zorgt ervoor dat de tijd die nodig is om naar een tweede set volumes over te stappen vermindert, terwijl het z/OS systeem actief blijft samen met de actieve applicaties \autocite{dhondy2005gdps}.
Deze functionaliteit implementeren verkort dan ook de RTO bij falen op opslagniveau, onder andere doordat er niet handmatig moet gewisseld worden van volumes.
Om over te schakelen naar site C kan de RTO ook verminderd worden doormiddel van geautomatiseerde overschakelingsmechanismen.
Frequente consistency groups voor de asynchrone site C zorgen er ook voor dat de RPO kan verkleind worden.
Als er elke 30 seconden een consistency group wordt gemaakt, bijvoorbeeld, dan is de RPO maximaal 30 seconden.
Daarom kan het nuttig zijn om vaak consistency groups te maken, omdat de RPO zo verlaagd kan worden.
Echter moet bij het kiezen van de frequentie van consistency groups ook gedacht worden aan de impact op het netwerk.
Een kleinere intervalwaarde zal er dus voor zorgen dat de sessie vaker probeert consistency groups te maken, wat ook de verwerkingssnelheid en belasting op netwerkverkeer op de opslagsystemen kan verhogen \autocite{ibm_2026_ibm_sessions_consistency_groups}.
De intervalwaarde kan ook ingesteld worden op nul, en zal er voor zorgen dat consistency groups continu worden aangemaakt.
Dit betekent dat, zodra de vorige consistency group is opgesteld, er een nieuwe consistency group wordt aangemaakt \autocite{ibm_2023_ibm_mirror_global_consistency_groups}.
De frequentie van consistency groups wordt best afgestemd op de crucaliteit van applicaties.
Voor kritieke systemen, zoals betalingsverwerking, is bijvoorbeeld een lagere interval aanbevolen.
Wanneer er moet worden overschakeld naar site C, moet ook rekening gehouden worden met de consistency groups.
Site C kan pas operationeel worden zodra alle consistency groups volledig zijn gerepliceerd.
Automatisatie, bij het pauzeren van updates die verstuurd worden naar de herstelcite, kan helpen om de herstelsite (site C) zo snel mogelijk operationeel te krijgen.
Het is mogelijk dat er inconsistenties onstaan bij het overschakelen naar de herstelcite.
Hierbij kan Controlled Log Catch-Up gebruikt worden om deze inconsistenties te kunnen corrigeren.
Bij het overschakelen naar site C, zullen applicaties herstart moeten worden, afhankelijk van de opstelling van het systeem.
Dit proces automatiseren kan er ook voor zorgen dat de RTO minder wordt.

De replicatievertraging van site B naar site C is idealiter zo laag mogelijk.
Daarom moet ook gekeken worden naar de replicatie op netwerkniveau.
Om de vertraging te minimaliseren kunnen compressietechnieken helpen om de hoeveelheid data die naar site C wordt verzonden te beperken.
Hiervoor zijn directe implementaties beschikbaar op hardware-niveau en op software-niveau \autocite{clayton_2016_ds8000}.
Volgens \textcite{ecklund_2024_unlocking} biedt de implementatie op hardware-niveau van de IBM DS8000 (10e generatie) een compressieverhouding van 2:1.
Implementaties op software-niveau, daarentegen, kunnen hogere compressieverhoudingen bereiken, maar zijn doorgaans minder snel dan implementaties op hardware-niveau.
Hoewel software-implementaties meer compressie bieden, vragen ze ook meer verwerkingstijd, wat de totale netwerkvertraging kan vergroten.
Hierdoor is de hardware-implementatie de voorkeursoptie omdat het een lagere impact heeft op de prestaties.
Bovendien zit de compressie ingebouwd in de opslagsystemen van IBM, zoals de DS8000 (10e generatie).
Compressie is niet enkel een voordeel bij een hybrideopstelling waarbij GDPS Metro Mirror en GDPS Global Mirror worden gecombineerd, maar ook bij alle besproken configuraties ongeacht de gekozen implementatie (hardware of software).

Om replicatieproblemen te kunnen detecteren en er snel op te kunnen reageren is het monitoren van replicatieprestaties en meldingen creëren indien nodig noodzakelijk.
Een mogelijkse tool die kan gebruikt worden is \textbf{IBM Z Operational Log and Data Analytics} \autocite{ibm_2026_z_operational_log_data_analytics_overview}.
Het kan gebruikt worden om operationele data, logbestanden en meetbare waarden in verband met prestaties te analyseren.
IBM biedt ook andere monitoring tools zoals \textbf{IBM Z IntelliMagic Vision} \autocite{ibm_2026_zos_virtual_tape_performance_management_z_intellimagic_vision}.
Dit product laat toe om netwerklatentie te monitoren, inclusief het bijhouden van hoeveel datareplicatie achterloopt en hoe lang het duurt voordat kopieën zijn verwerkt, wat relevant is voor disaster recovery en het waarborgen van RPO-doelstellingen.
Er is ook nog \textbf{IBM Z Performance and Capacity Analytics}, waarmee de performantie van van een systeem kan beheerd worden \autocite{ibm_2026_ibm_z_performance_and_capacity_analytics}.
Het is ook mogelijk om er anomalieën in log- en metrische data te detecteren, wat indirect kan helpen om replicatieproblemen te identificeren.

\section{\IfLanguageName{dutch}{Risicoanalyse en Faalscenario's Asynchrone Replicatie in Db2 voor z/OS}{Risk Analysis and Fail Scenarios}}%
\label{sec:risico-analyse-asynchrone-log-replicatie}

In het vorige hoofdstuk zijn verschillende mainframe-opstellingen besproken, inclusief hun werking, voordelen, beperkingen en mogelijke faalscenario's waarbij het voornamelijk ging over opslaggebaseerde replicatie-opstellingen zoals GDPS Metro Mirror en GDPS Global Mirror.
Dit hoofdstuk richt zich op log-gebaseerde replicatie in Db2 voor z/OS, met name Q Replication.
Hierbij wordt gekeken naar de risico's die gepaard gaan met asynchrone replicatie op transactieniveau, zoals replicatievertraging, incomplete LUW's en inconsistenties in logreplicatie.
Deze risico's worden gekoppeld aan kritieke componenten zoals log streams en coupling facilities, en hun impact op RPO, RTO en transactionele consistentie wordt geanalyseerd.

\subsection{\IfLanguageName{dutch}{Replicatievertraging}{Replication Latency}}
\label{sec:asynchrone-log-replicatie-vertraging}

Replicatievertraging is een risico bij asynchrone logreplicatie, zoals Q Replication in Db2 voor z/OS door de asynchrone aard van de technologie.
In tegenstelling tot een synchrone manier van werken, waarbij op bevestiging wordt gewacht, werkt Q Replication asynchroon en wordt dus niet op bevestiging gewacht.
Hierbij wordt er met MQ queues gewerkt die systematisch worden aangemaakt en verstuurd naar het Q Apply programma die de informatie in de queues terug omzet naar transacties \autocite{ibm_2026_infosphere_q_replication}.
Replicatievertraging verwijst specifiek naar de tijdsduur tussen het committen van een transactie op de primaire Db2 databank en het succesvol toepassen van die transactie op de secundaire databank.
Deze vertraging is inherent aan asynchrone replicatie en wordt beïnvloed door factoren zoals het netwerk, MQ of Db2.

Volgens \textcite{sharkey_2018_ibm_performance_q_replication} is de impact op de vertragingstijd bij het repliceren met Q Replication reëel.
De afstand tussen de sites en het aantal berichten dat verstuurd wordt, de zogenaamde depth of diepte van een queue, heeft hierbij een invloed op de vertragingstijd.
De bandbreedte van de connectie tussen de twee sites kan ook onvoldoende zijn en kan er voor zorgen dat een hoog volume aan logrecords trager wordt verwerkt.
IBM MQ werkt synchroon op pakketniveau, omdat TCP/IP, het onderliggende protocol, volgens \textcite{borman_2014_tcp} wacht op bevestigingen (ACKs) voor elk verzonden pakket. 
Als een pakket verloren gaat, wordt het herverzonden, wat de doorvoersnelheid van berichten kan vertragen.
Q Replication daarentegen is asynchroon op transactieniveau.
De primaire database wacht namelijk niet op bevestiging van Q Apply (de toepassing op de secundaire database) voordat nieuwe transacties worden verwerkt. 
Hierdoor blijft de primaire database operationeel, terwijl de replicatie op de achtergrond plaatsvindt. 
De TCP/IP herverzending kan dus zorgen voor een vertraging van de doorvoersnelheid.
Er kan ook vertraging ontstaan bij het verwerken van MQ queues.
Als er veel onverwerkte berichten in de MQ queue staan, moet Q Apply alle berichten sequentieel verwerken.
Hoe dieper de queue, hoe langer het dus duurt voordat een bericht wordt verwerkt.
Hierbij is de impact op vertraging evident, want berichten kunnen pas verwerkt worden nadat de vorige berichten zijn verwerkt.
Diepe queues kunnen er ook voor zorgen dat disk en CPU gebruik verhogen waardoor de verwerking ook vertraagt \autocite{sharkey_2018_ibm_performance_q_replication}.

Db2 specifieke factoren kunnen er ook voor zorgen dat de vertragingstijd groter wordt.
Grote logrecords, bijvoorbeeld bij een operatie op een grote tabel waarbij alle waarden worden bewerkt met een UPDATE operatie, kunnen de overhead verhogen.
Grotere logrecords vereisen meer bandbreedte en verwerkingstijd in Q Capture en MQ \autocite{ibm_2004_ibm_replication_event_publishing_guide}.
Aanvullend, kunnen applicaties die kleine frequente commits doen meer CPU-verbruik veroorzaken \autocite{ibm_2025_infosphere_batch_frequency_param}.
Naast grote logrecords en kleine frequente commits, is er ook nog een mogelijke vertraging wanneer de log buffers van Db2 vol zijn.
Als dit gebeurt, moet Db2 immers wachten tot de log buffers naar de active log zijn geschreven \autocite{ibm_2026_how_db2_writes_active_logs}.

De performantie van de site waarnaar de queues worden verstuurd heeft ook een invloed op de mogelijke vertragingstijd.
De secundaire site moet namelijk met Q Apply transacties toepassen op de secundaire database.
Hierbij is de snelheid afhankelijk van zaken zoals CPU en nodige operaties naar opslagapparatuur (I/O).
Tabellen met veel indexes kunnen ook het toepassen van transacties vertragen, en constraints zoals bijvoorbeeld foreign keys kunnen ervoor zorgen dat er extra validatie nodig is.
Bij Q Apply kan het toevoegen van nieuwe indexen een invloed hebben op de performantie en kan het zorgen voor meer nodige controles \autocite{ballard_2016_understanding_q_replication_high_availability}.

\subsubsection{\IfLanguageName{dutch}{Meten Replicatievertraging}{Measuring Replication Latency}}
\label{sec:asynchrone-log-replicatie-meten-replicatievertraging}

Om replicatievertraging te meten en te monitoren kunnen de volgende Key Performance Indicators (KPI's) en tools gebruikt worden.
Voor een duidelijk begrip van de tabel met de KPI's is een uitleg van de verschillende besproken tools en commando's vereist.
Om te beginnen is er de Q Replication Alert Monitor  die kan worden gestart met het \textbf{asnmon}-commando via de UNIX System Services (USS) op z/OS \autocite{ibm_2026_asnmon}.
Deze monitoring tool gaat de status van replicatie na en kan melden maken wanneer er bepaalde zaken gebeuren in de replicatieomgeving \autocite{ibm_2026_q_replication_alert_monitor}.
Het kan bijvoorbeeld een melding genereren wanneer het Q Capture programma een vooraf gedifinieerde tijd heeft overschreden bij het doorgeven van de transactie aan het Q Apply programma \autocite{ibm_2026_alert_conditions_replication_alert_monitor}.
Daarbovenop laat de tool ook toe om een melding te creëren wanneer het Q Apply programma een bepaalde vooraf gedifinieerde tijd heeft overschreden bij het toepassen van de transacties op de secundaire databank.

IBM biedt een tool genaamd IMB MQ Explorer die draait op systemen met Windows of Linux waarmee MQ objecten kunnen worden gemonitord op verschillende platformen waaronder z/OS \autocite{ibm_2026_ibm_mq_explorer}. 
Het laat ook toe om de diepte van MQ queues te bekijken \autocite{ibm_2026_displaying_queue_channel_monitoring_data}.

Om de vertragingstijd op netwerkniveau te meten zijn er verschillende mogelijkheden waaronder de volgende commando's en tools.
\begin{itemize}
    \item Het \textbf{ping}-commando op USS voor z/OS kan gebruikt worden om de omlooptijd van pakketten tussen twee systemen te meten \autocite{ibm_2026_the_uss_zos_ping_command}. Het toont de vertragingstijd in milliseconden. Er zijn wel beperkingen, want het ping-commando is niet representatief voor TCP/IP-verkeer zoals bij MQ. Het toont ook geen verdere informatie over de routepaden in het netwerk en mogelijke knelpunten.
    \item Het \textbf{traceroute}-commando toont het pad dat een pakket neemt naar een gevraagde bestemming met daarbij het aantal \textbf{hops} of routers dat er moeten worden gepasseerd \autocite{ibm_2026_using_the_traceroute_command}. Het commando kan uitgevoerd worden via TSO met \textbf{TRACERTE} of via een USS voor z/OS shell met \textbf{traceroute}.
    \item IBM Z NetView biedt verschillende functionaliteiten om IBM Z netwerken beschikbaar te houden \autocite{ibm_2026_ibm_z_netview}. Het kan onder andere netwerkvertraging monitoren tussen z/OS-systemen, inclusief TCP/IP-verkeer.
\end{itemize} 

In Db2 is er de DSN1LOGP utility die kan gebruikt worden om de inhoud van recovery logs te formatteren voor weergave \autocite{ibm_2026_db2_DSN1LOGP}. 
Er zijn twee mogelijke weergaven zoals een gedetailleerd rapport van alle individuele logs en een samenvattend rapport.
De Db2 Log Analysis Tool laat toe om veranderingen in tabellen te monitoren en hier automatisch rapporten van te maken \autocite{ibm_2025_ibm_log_analysis_tool}. 
Het kan gebruikt worden om ongewenste veranderingen in Db2 tabellen te isoleren.

\begin{table}[h]
\centering
\begin{tabularx}{\textwidth}{|p{2cm}|p{4cm}|p{4cm}|X|}
\hline
\textbf{Meetbare waarde} & \textbf{Beschrijving} & \textbf{Tool/Commando} & \textbf{Streefcijfer} \\
\hline
End-to-End Latency & Tijd tussen commit op primaire site en toepassen op secundaire site via de \textbf{QCAPTURE\_LATENCY} en \textbf{QAPPLY\_LATENCY} alert condities. & Q Replication Alert Monitor & Minder dan één seconde \\
\hline
MQ Queue Depth & Aantal onverwerkte berichten in de MQ queue. & MQ Explorer & Minder dan 100 berichten \\
\hline
Network Latency & Round-trip tijd (RTT) tussen primaire en secundaire site. & ping, traceroute, of IBM Z NetView & Minder dan 50 milliseconden lokaal, minder dan 200 milliseconden wereldwijd \\
\hline
Log Generation Rate & Aantal logrecords gegenereerd per seconde op de primaire site. & DSN1LOGP of Db2 Log Analysis Tool & Afhankelijk van workload \\
\hline
\end{tabularx}
\caption{Meetbare waarden voor replicatievertraging}
\label{tab:q_replication_metrics}
\end{table}

\subsubsection{\IfLanguageName{dutch}{Impact van Replicatievertraging}{Impact of Replication Latency}}
\label{sec:asynchrone-log-replicatie-impact-replicatievertraging}

De impact van replicatievertraging is direct gekoppeld aan RPO, RTO en operationele veerkracht.
Hoe groter de vertraging, hoe meer dataverlies optreedt bij een uitvalscenario.
Als het dataverlies te groot is, kan dit een risico vormen voor financiële instellingen omdat er niet voldaan wordt aan wettelijke eisen zoals DORA vanuit de Europese Unie.
Vertraagde failover, mogelijks veroorzaakt door het moeten inhalen van een achterstand (bijvoorbeeld met Controlled Log Catch-Up), betekent een langere downtime waardoor er zich financiële verliezen en reputatieschade kunnen voordoen.
Replicatievertraging of andere replicatiefouten kunnen instaan voor transacties die niet in de juiste volgorde zijn toegepast. 
Dit zorgt voor inconsistente data en kan foutieve saldi, dubbele transacties of ontbrekende data met gevolg hebben bij systemen van financiële instellingen.
Als er transacties moeten worden ingehaald door replicatievertraging, kan dit ook een impact hebben op de prestaties van applicaties omdat er in korte tijd veel CPU- en I/O-belasting veroorzaakt wordt.
Dit heeft op zijn beurt mogelijkse impact op cruciale batchprocessen.
Als de replicatievertraging te groot is, kan dit betekenen dat er bepaalde opgelegde wetten niet worden nageleefd.
Voor financiële instellingen kan dit boetes tot gevolg hebben.

\subsubsection{\IfLanguageName{dutch}{Mitigatiestrategieën voor Replicatievertraging}{Mitigation Strategies for Replication Latency}}
\label{sec:asynchrone-log-replicatie-mitigatiestrategie-replicatievertraging}

Er zijn verschillende manieren om replicatievertraging te minimaliseren.
Een van de basis mitigatiestrategieën is het optimaliseren van het netwerk waarop replicatie wordt gedaan.
Het gebruik van afgescheiden netwerken met dark fibre is sterk aangeraden.
Replicatieverkeer prioriteren boven andere applicaties kan er ook voor zorgen dat er minder replicatievertraging is.
Daarnaast kan de voorheen besproken hardware- en software-compressie helpen om het bandbreedtegebruik te verminderen zodat er efficiënter wordt gerepliceerd.
Q Capture leest logrecords uit de Db2 log en groepeert deze in batches voordat ze als een gezamenlijk MQ-bericht worden verzonden \autocite{ibm_2025_infosphere_batch_frequency_param}.
Het verlagen van de batchgrootte kan voor snellere verwerking zorgen, maar veroorzaakt wel meer netwerkoverhead.
Q Capture laat ook toe om parallele queues te maken \autocite{ibm_2026_infosphere_parallel_send_queues}.
Parallele queues vermijden ophoping op één MQ queue.
Door het parallele versturen kan de vertragingstijd verminderd worden en de transactiedoorvoer worden verbeterd.
Meerdere queues kunnen tot meer CPU en geheugen gebruik leiden.
Indien het risico acceptabel is, kan er ook voor gekozen worden om berichten niet persistent te maken.
Dit betekent dat berichten niet naar disk worden gekopieerd wanneer er zich een fout voordoet \autocite{ibm_2025_ensuring_messages_not_lost_mq}.
Berichten die niet persistent zijn gaan verloren wanneer het systeem of een queue manager faalt.
Voor IBM MQ zijn grotere buffer pools algemeen beter voor performantie \textcite{ibm_2026_ibm_defining_buffer_pools}.
Grotere buffer pools kunnen ervoor zorgen dat er minder I/O-operaties moeten gedaan worden, wat resulteert in minder vertraging.
Het vergroten van buffer pools zorgt voor meer geheugengebruik.

\subsection{\IfLanguageName{dutch}{Incomplete Logical Units of Work}{Incomplete Logical Units of Work}}
\label{sec:asynchrone-log-replicatie-incomplete-logical-units-of-work}

Een ander mogelijk risico bij logreplicatie zijn incomplete Logical Units of Work (LUW's).
Een LUW wordt doorgaans beëindigt door een commit, een volledige rollback van een transactie of het einde van een applicatie \autocite{ibm_2026_db2_unit_of_work}.
Als de veranderingen nog niet committed zijn, kunnen andere applicaties deze veranderingen nog niet zien.
Het start- en stoppunt van een LUW definieert consistentie punten.
\textcite{ibm_2026_db2_unit_of_work} geeft het volgend voorbeeld.
Een transactie binnen een bank kan inhouden dat er een bepaalde som aan geld wordt overgeplaatst van een rekening naar een andere rekening.
Bij een dergelijke transactie moet eerst de som geld afgetrokken worden van de ene rekening en dan toegevoegd worden aan de andere rekening.
Nadat het geld afgetrokken is van de ene rekening, wordt de data inconsistent want het geld moet nog worden toegevoegd op de tweede bankrekening.
Alleen wanneer het geld afgetrokken is en correct toegevoegd op de desbetreffende rekeningen, kan de consistentie van de data weer bevestigd worden.
Wanneer beide transacties, in dit geval, voltooid zijn, kan een commit gedaan worden om de transactie af te ronden waarna de veranderingen ook zichtbaar worden voor alle applicaties.
Als er een fout gebeurd tijdens het uitvoeren van de transactie, zal Db2 gewoonlijk de wijzigingen terug draaien en daarmee de data terugzetten naar de start van de LUW \autocite{ibm_2026_how_rolls_back}.
Het probleem ontstaat wanneer de LUW niet kan worden voltooid en, tijdens een rollback, er een fout ontstaat waardoor de data niet terug gezet is.
Hierdoor kan de data inconsistent zijn.
Bij het voorbeeld van \textcite{ibm_2026_db2_unit_of_work} zou het geld afgetrokken zijn van de eerste rekening maar niet toegevoegd zijn aan de tweede rekening.
Als een rollback dan niet lukt, is de data inconsistent.
Het kan gebeuren dat een asynchrone replicatieomgeving zoals Q Replication LUW's daardoor onvolledig repliceerd.

\subsubsection{\IfLanguageName{dutch}{Oorzaken van Incomplete LUW's}{Causes of Incomplete Logical Units of Work}}
\label{sec:asynchrone-log-replicatie-oorzaken-incomplete-logical-units-of-work}

Incomplete LUW's kunnen optreden door verschillende storingen in het replicatieproces van Q Replication.
Hieronder worden de hoofdoorzaken besproken, samen met hun impact op de dataconsistentie.

Een mogelijke oorzaak van incomplete LUW's is een crash van het Q Apply programma tijdens het toepassen van een transactie op de secundaire databank.
Q Apply is verantwoordelijke voor het toepassen van de logrecords, die door Q Capture zijn verzameld en via MQ zijn verzonden, op de secundaire databank.
Het MQ-bericht dat Q Apply ontvangt bevat een volledige LUW.
Q Apply begint, na het ontvangen van het bericht, met het toepassen van de individuele operaties binnen die LUW in de juiste volgorde.
Het probleem ontstaat wanneer Q Apply crasht tijdens dit proces, bijvoorbeeld door een systeemfout.
Hierdoor kan het voorkomen dat er enkel een deel van de operaties binnen de LUW toegepast is toegepast, met inconsistente data als gevolg.

Een andere oorzaak van incomplete LUW's zijn gaten in de logbestanden op de primaire site.
Db2 voor z/OS schrijft alle wijzingingen in de databank naar de logbuffer die uiteindelijk worden bewaard in de active log en later naar de archive log worden gekopieerd.
Als er een schrijffout optreedt, door bijvoorbeeld het falen van een opslagapparaat, kan het zijn dat een of meerdere logrecords ontbreken in de active of archive log \autocite{broadcom_2026_db2_backup_recovery_handbook}.
Wanneer Q Capture de logrecords leest om LUW's te reconstrueren en naar de secundaire site verzenden, mist het de ontbrekende logrecords.
Hierdoor kan Q Capture een onvolledige LUW naar Q Apply verzenden.
Als Q Apply deze onvolledige LUW toepast op de secundaire database, onstaat er datainconsistentie.
Dit heeft directe gevolgen voor de RPO.
Het herstellen van een dergelijke situatie is complex en vereist het gebruik van backups en archive logs om de ontbrekende gegevens te reconstrueren.
Met nae als de backups niet actueel zijn, kan de RTO aanzienlijk verlengd worden.

In omgevingen met een hoge workload kan de MQ-queue die gebruikt wordt voor Q Replication overbelast worden.
Q Capture verzendt LUW's met MQ-berichten naar de secundaire site, waar Q Apply deze berichten verwerkt \autocite{ibm_2026_infosphere_q_replication}.
Als Q Capture sneller berichten verzendt dan Q Apply ze kan verwerken, groeit de queue tot een maximale grootte \autocite{ballard_2016_understanding_q_replication_high_availability}.
Dit kan leiden tot een langere replicatietijd, zoals voorheen gesproken, maar kan ook leiden tot ontbrekende LUW's als het systeem faalt wanneer bepaalde LUW's nog niet zijn toegepast door Q Apply.

\subsubsection{\IfLanguageName{dutch}{Impact van Incomplete LUW's}{Impact of Incomplete Logical Units of Work}}
\label{sec:asynchrone-log-replicatie-impact-incomplete-logical-units-of-work}

Een storing veroorzaakt door incomplete LUW's heeft meerdere gevolgen.
Ten eerste ontstaat er datainconsistentie, wat voor financiële instellingen kan leiden tot onjuiste saldi, dubbele uitgevoerde operaties, of ontbrekende gegevens.
Ten tweede heeft dit gevolgen voor de RPO en RTO.
Alle transacties die na de onvolledige LUW zijn gemmit op de primaire site, zijn mogelijk verloren als de primaire site uitvalt voordat de replicatie is hersteld.
Bovendien vereist het herstellen van de consistentie vaak handmatige interventie of het uitvoeren van Controlled Log Catch-Up (CLC), wat de RTO aanzienlijk kan verlengen.

\subsubsection{\IfLanguageName{dutch}{Mitigatiestrategieën voor Incomplete LUW's}{Mitigation Strategies for Incomplete Logical Units of Work}}
\label{sec:asynchrone-log-replicatie-mitigatie-incomplete-logical-units-of-work}

De risico's die gepaard gaan met incomplete LUW's in een asynchrone replicatieomgeving zoals Q Replication kunnen aanzienlijk worden gemitigeerd door een combinatie van preventieve en correctieve maatregelen. 
Voor financiële instellingen, waar dataconsistentie en operationele veerkracht van het grootste belang zijn, is het implementeren van deze strategieën essentieel om te voldoen aan strenge eisen zoals die gesteld worden door DORA, PCI-DSS en GDPR. 
Hieronder worden de meest effectieve mitigatiestrategieën besproken.

Preventieve maatregelen zijn gericht op het voorkomen van situaties die kunnen leiden tot incomplete LUW's. 
Een van de meest effectieve strategieën is het gebruik van frequente consistency groups in combinatie met GDPS Global Mirror. 
Door elke 30 seconden een consistency group aan te maken, wordt de RPO beperkt tot maximaal 30 seconden. 
Dit zorgt ervoor dat, zelfs bij een storing, het dataverlies beperkt blijft tot de transacties die plaatsvonden tussen twee opeenvolgende consistency groups. 
Hoewel deze aanpak de netwerkbelasting verhoogt, is de winst in termen van databescherming aanzienlijk, met name voor kritieke systemen \autocite{gomes_2025_redbooks}.

Redundante netwerkverbindingen kunnen worden geïmplementeerd om het risico op netwerkstoringen te minimaliseren. 
Door twee of meer onafhankelijke netwerkpaden te gebruiken voor replicatie, kan het systeem overschakelen naar een backup-verbinding als de primaire verbinding uitvalt. 
Dit vermindert de kans dat MQ-berichten verloren gaan of onvolledig worden ontvangen \autocite{ballard_2016_understanding_q_replication_high_availability}.
Tot slot kan de batchgrootte in Q Capture worden geoptimaliseerd. 
Q Capture leest logrecords uit de Db2-log en groepeert deze in batches voordat ze als MQ-berichten worden verzonden. 
Een kleinere batchgrootte zorgt voor snellere verwerking, maar veroorzaakt meer netwerkoverhead. 
Omgekeerd vermindert een grotere batchgrootte de overhead, maar kan de verwerking vertragen. 
Het is daarom aan te raden om de batchgrootte af te stemmen op de specifieke workload en netwerkomstandigheden \autocite{ibm_2025_infosphere_batch_frequency_param}.

Indien incomplete LUW's toch optreden, zijn er verschillende correctieve maatregelen beschikbaar om de dataconsistentie te herstellen. 
De RECOVER utility in Db2 laat toe om data herstellen naar een huidige of vorige staat door het kopiëren van en toepassen van logrecords \autocite{ibm_2025_db2_recover}.
Als een LUW onvolledig is toegepast, kan de RECOVERY utility deze terugdraaien of opnieuw toepassen om de consistentie te herstellen \autocite{ibm_2021_db2_clc}.
Hoewel deze aanpak tijdrovend kan zijn, is het vaak noodzakelijk om de dataconsistentie te waarborgen \autocite{ibm_2026_how_rolls_back}.

In het ergste geval kan de secundaire database hersteld worden vanaf een backup. 
Dit is echter een laatste redmiddel, omdat alle transacties die sinds de laatste backup zijn uitgevoerd, verloren zullen gaan. 
Het herstelproces omvat het herstellen van de database vanaf de backup, het toepassen van ontbrekende transacties met behulp van archive logs, en het opnieuw synchroniseren met de primaire database. 
Deze aanpak heeft een aanzienlijke impact op zowel de RPO als de RTO en moet daarom alleen worden overwogen wanneer andere herstelmethoden niet beschikbaar zijn \autocite{broadcom_2026_db2_backup_recovery_handbook}.

Tot slot is regelmatige monitoring en testing essentieel. 
Door replicatieprestaties te monitoren met tools zoals IBM Z NetView en Q Replication Alert Monitor, kunnen potentiële problemen tijdig worden gedetecteerd. 
Daarnaast is het raadzaam om failover-scenario's minstens één keer per jaar te testen om te waarborgen dat de RPO en RTO binnen de gestelde doelen blijven.

\subsection{\IfLanguageName{dutch}{Inconsistenties Volgorde Logreplicatie}{Inconsistencies Order Log Replication}}
\label{sec:asynchrone-log-replicatie-inconsistenties-volgorde-logreplicatie}

Bij asynchrone replicatie, zoals Q Replication, worden transacties niet onmiddelijk gesynchroniseerd tussen een primaire en secundaire site.
Dit introduceert een vertraging, zoals eerder besproken, tussen het moment waarop een transactie wordt gecommit op de primaire databank en het moment waarop de transactie wordt toegepast op de secundaire databank.
Het is essentieel dat de volgorde van de transacties wordt behouden ondanks deze replicatievertraging.
Met name in financiële instellingen zijn transacties vaak afhankelijk van elkaar.
Bij een overschrijving, bijvoorbeeld, moet eerst het geld worden afgenomen van de ene rekening voordat het wordt toegevoegd aan de andere rekening.
Als deze volgorde niet work gerespecteerd kan dit leiden tot zaken zoals negatieve saldi en dubbele betalingen.
Bovendien mag asynchrone replicatie de ACID-principes (Atomicity, Consistency, Isolationen Durability) niet overschrijden. 
Ook moeten na een storing de transacties in de juiste volgorde kunnen worden hersteld om RPO en RTO te waarborgen.
Een verkeerde volgorde kan leiden tot dataverlies of langere hersteltijden.

\subsubsection{\IfLanguageName{dutch}{Logrecordverwerking in Db2}{Logrecord Processing in Db2}}
\label{sec:asynchrone-log-replicatie-logrecordverwerking-db2}

Q Capture leest de logrecords uit de Db2 logs en groepeert deze in LUW's.
Elke LUW wordt omgezet in een MQ-bericht en verzonden naar de secundaire site via IBM MQ.
Belangrijk om te vermelden is dat Q Capture aan elk MQ-bericht metadata toevoegt.
Deze metadata bevat een sequence nummer, een uniek nummer dat de volgorde van het bericht binnen de LUW aangeeft, een transaction ID waarmee de LUW uniek geïdentificeerd kan worden, en het tijdstip waarop de transactie werd gecommit op de primaire site \autocite{ibm_2025_structure_delimited_messages_event_publishing}.
Q Apply ontvangt de MQ-berichten op de secundaire site en past de wijzigingen toe op de secundaire databank.
Het is cruciaal dat Q Apply de berichten verwerkt in de juiste volgorde om de transactionele consistentie te waarborgen.
Standaard verwerkt Q Apply berichten sequentieel, maar parallele threads kunnen ook gebruikt worden voor verschillende LUW's \autocite{ibm_2026_q_apply_program}.

IBM MQ biedt twee mechanismen om de volgorde van berichten te behouden.
Het sequence nummer zorgt ervoor dat berichten in volgorde kunnen worden geleverd.
MQ kan berichten ook groeperen in een message group.
Dit laat toe om berichten te markeren als gerelateerd aan elkaar.
Daarbovenop kan er ook een logische volgorde aan de groep toegepast worden \autocite{ibm_2026_message_group}.
MQ waarborgt verzekert dat de berichten in deze volgorde zullen worden geleverd.

\subsubsection{\IfLanguageName{dutch}{Oorzaken en Mitigatiestrategieën Inconsistensies Volgorde}{Causes of Inconsistencies in Order}}
\label{sec:asynchrone-log-replicatie-oorzaken-inconsistenties-volgorde}

IBM MQ biedt mechanismen om de volgorde te bewaren, maar verkeerde configuraties kunnen leiden tot inconsistenties.
Berichten kunnen in willekeurige volgorde toekomen en daarom is het inschakelen van message sequencing noodzakelijk om de juiste volgorde te bewaren \autocite{ibm_2025_message_sequencing}.
Als message sequencing niet is ingeschakeld, zijn berichten niet gegarandeerd in de correcte volgorde.
Schakel daarom message sequencing in voor alle queues die worden gebruikt door Q Replication.

Een andere mogelijke oorzaak door misconfiguratie heeft te maken met message groups.
Als berichten van één LUW niet in dezelfde message group worden geplaatst, kan Q Apply ze in verkeerde volgorde verwerken.
Dit kan gebeuren als de \textbf{MESSAGE\_FORMAT} in de \textbf{IBMQREP\_SENDQUEUES} tabel niet op C (compact) staat \autocite{bibid}{ibm_2026_ibmqrep_sendqueues}.
Ook kan het zijn dat de \textbf{MSG\_CONTENT\_TYPE} verkeerd is ingesteld.
Als \textbf{MSG\_CONTENT\_TYPE} namelijk op R staat, worden berichten per operatie (zoals INSERT, UPDATE en DELETE) verzonden.
Hierdoor worden berichten door Q Capture niet als één transactiegroep herkent.

Fouten in de configuratie of uitvoering van Q Apply kunnen ook leiden tot inconsistenties met betrekking tot de volgorde van transacties.
Q Apply kan meerdere threads gebruiken om LUW's toe te passen.
Als afhankelijke LUW's parallele worden verwerkt, kan dit leiden tot inconsistenties.
Om dit te vermijden, moeten afhankelijke tabellen in één replication queue map worden geplaatst, zodat Q Apply de volgorde kan garanderen.

